\subsection{A Fixed-Lag Estimate of the Family at an Alarm}
\label{sec:alarm-family}

Once an alarm has been raised, a later PMU report can clarify which
family was active at that alarm, even if the source has since cleared.
Let $b$ be the report at the alarm and $b_-$ the preceding alarm, or
$-1$ at startup. For each retained history $h$, let $u_h$ and $v_h$
be the onset and clearance report of its latest intervention, with
$v_h=\infty$ while it remains active. The histories consistent with a
fresh intervention active at $b$ satisfy $b_-<u_h\le b<v_h$.
At publication report $k=b+L$, assign family $f$ the mass
\begin{equation}
 q_f(b\mid k)=\sum_{h:\, f(h)=f}\pi_{h,k}\,
        \mathbf{1}\{b_-<u_h\le b<v_h\}.
 \label{eq:alarm-family-mass}
\end{equation}
Here $\pi_{h,k}$ includes all reports through $k$. Two candidate load
buses can therefore support the same diagnosis even when neither bus
alone has the highest weight. A load restored between $b$ and $k$ still
supports a load diagnosis at $b$.

Write $Q=\sum_f q_f$. Conditional on the retained histories satisfying
the alarm condition, the expected unit loss from declaring family $g$
is $(Q-q_g)/Q$. Thus $\arg\max_g q_g$ minimizes family error for this
local model when $Q>0$. With no retained support the estimator abstains.
This fixed-lag estimate changes the family attached to an existing alarm.
The alarm clock, its threshold and the pruning approximation are inherited
from the bank. The lag fixes the information available to the classifier
and is charged in full to its publication time.
